\chapter{Modelo de datos del perfil enriquecido}
\label{annex:modelo-datos}

Este anexo detalla el documento de perfil que la fase de consolidación produce para cada investigador y que constituye la unidad de trabajo de la indexación y la evaluación.

\begin{table}[htbp]
    \centering
    \caption{Campos principales del documento de perfil enriquecido.}
    \label{tab:perfil-campos}
    \begin{tabularx}{\textwidth}{l X}
        \toprule
        \rowcolor{upmblue!8} \colhead{Campo} & \colhead{Contenido} \\
        \midrule
        Identificador & Identificador público del investigador en el portal. \\
        Filiación & Posiciones ocupadas y unidades organizativas (escuela, departamento). \\
        Temas & Lista normalizada de temas de investigación declarados (variantes y duplicados unificados). \\
        Publicaciones & Producción científica con título, año y resumen depurado. \\
        Proyectos supervisados & Trabajos dirigidos recientemente, con título, resumen y palabras clave (base de la verdad de referencia del Capítulo~\ref{ch:evaluacion}). \\
        Dominios & Áreas de especialización extraídas automáticamente, acumulativas entre ejecuciones. \\
        Biografía & Síntesis en prosa de la trayectoria, generada incrementalmente. \\
        Metadatos & Fechas de actividad (primer y último año), recuentos y trazas de procesamiento. \\
        \bottomrule
    \end{tabularx}
\end{table}

Antes de consolidarse en ese documento, la evidencia se modela como un \emph{property graph} (Sección~\ref{sec:met-grafo}).
La Figura~\ref{fig:grafo-esquema} muestra su esquema: un nodo central \texttt{Pdi} (investigador) conectado con los cuatro tipos de entidad que describen su actividad, más una relación de coautoría entre investigadores.

\begin{figure}[htbp]
    \centering
    \begin{tikzpicture}[
        font=\small,
        nodo/.style={circle, draw=upmblue!80!black, fill=upmblue!10, thick, minimum size=16mm, align=center, inner sep=1.5pt, font=\footnotesize\sffamily},
        central/.style={circle, draw=upmblue!85!black, fill=upmgold!35, very thick, minimum size=18mm, align=center, font=\bfseries\sffamily},
        rel/.style={-{Stealth[length=2.4mm]}, semithick, draw=upmblue!70!black},
        rellab/.style={font=\scriptsize\sffamily, text=upmblue!85!black, fill=white, inner sep=1.5pt}
    ]
        \node[central] (pdi)  at ( 0, 0)    {Pdi};
        \node[nodo]    (pub)  at (-4, 2.4)  {Publi-\\cation};
        \node[nodo]    (tft)  at ( 4, 2.4)  {LastYear-\\Project};
        \node[nodo]    (pos)  at (-4,-2.4)  {Position};
        \node[nodo]    (subj) at ( 4,-2.4)  {Subject};

        \draw[rel] (pdi) -- node[rellab,sloped]{AUTHORSHIP} (pub);
        \draw[rel] (pdi) -- node[rellab,sloped]{TUTORING}   (tft);
        \draw[rel] (pdi) -- node[rellab,sloped]{WORKS}      (pos);
        \draw[rel] (pdi) -- node[rellab,sloped]{TEACHING}   (subj);
        \draw[rel] (pdi) to[out=120,in=60,looseness=9] node[rellab,above]{COAUTHORED} (pdi);
    \end{tikzpicture}
    \caption{Esquema del grafo de propiedades: tipos de nodo y relaciones. Todas las relaciones del modelo parten del investigador (\texttt{Pdi}); \texttt{COAUTHORED} es una arista derivada que conecta investigadores entre sí.}
    \label{fig:grafo-esquema}
\end{figure}

La Tabla~\ref{tab:grafo-nodos} detalla la información que guarda cada nodo y la Tabla~\ref{tab:grafo-aristas} la de cada relación, conforme a los esquemas Pydantic del proyecto.

\begin{table}[htbp]
    \centering
    \caption{Tipos de nodo del grafo y su contenido.}
    \label{tab:grafo-nodos}
    \begin{tabularx}{\textwidth}{l l X}
        \toprule
        \rowcolor{upmblue!8} \colhead{Nodo} & \colhead{Identificador} & \colhead{Campos principales} \\
        \midrule
        \texttt{Pdi} & \texttt{pdi\_NNNNNN} & Nombre completo (y nombre/apellidos), correo, identificadores externos (ORCID, Scopus\ldots) y temas de investigación declarados en el portal. \\
        \texttt{Publication} & \texttt{pub\_N} & Título, descripción depurada, año, tipo (artículo, conferencia\ldots), fuente, agencia de validación, autores, palabras clave, e identificador y URL de acceso abierto. \\
        \texttt{LastYearProject} & \texttt{tft\_N} & Título, descripción, año, tipo, autores, palabras clave e identificador de acceso abierto. \\
        \texttt{Position} & \texttt{loc\_\ldots} & Nombre y tipo de entidad (escuela, departamento, grupo de investigación\ldots) e identificador de origen en el portal. \\
        \texttt{Subject} & \texttt{subj\_\ldots} & Nombre de la asignatura y tipo de docencia. \\
        \bottomrule
    \end{tabularx}
\end{table}

\begin{table}[htbp]
    \centering
    \caption{Relaciones del grafo: extremos y propiedades. \texttt{COAUTHORED} es una arista derivada (no proviene de un esquema Pydantic).}
    \label{tab:grafo-aristas}
    \begin{tabularx}{\textwidth}{l l X}
        \toprule
        \rowcolor{upmblue!8} \colhead{Relación} & \colhead{Origen $\rightarrow$ Destino} & \colhead{Propiedades} \\
        \midrule
        \texttt{AUTHORSHIP} & \texttt{Pdi} $\rightarrow$ \texttt{Publication} & Autoría de una publicación (sin propiedades adicionales). \\
        \texttt{TUTORING} & \texttt{Pdi} $\rightarrow$ \texttt{LastYearProject} & Supervisión de un TFG/TFM (sin propiedades adicionales). \\
        \texttt{WORKS} & \texttt{Pdi} $\rightarrow$ \texttt{Position} & Pertenencia o rol en una unidad: \texttt{role} y \texttt{since}. \\
        \texttt{TEACHING} & \texttt{Pdi} $\rightarrow$ \texttt{Subject} & Docencia de una asignatura: \texttt{start\_date} y \texttt{end\_date}. \\
        \texttt{COAUTHORED} & \texttt{Pdi} $\rightarrow$ \texttt{Pdi} & Coautoría entre investigadores: \texttt{num} (número de publicaciones en común) y \texttt{first\_year}/\texttt{last\_year} (primer y último año de colaboración). \\
        \bottomrule
    \end{tabularx}
\end{table}

\pendiente{Añadir un ejemplo real abreviado y anonimizado de un documento de perfil enriquecido (JSON validado por Pydantic).}

\chapter{Prompts de los agentes}
\label{annex:prompts}

Este anexo recoge las instrucciones (\emph{prompts}) de los tres agentes basados en LLM que intervienen en el sistema, con su contrato de salida estructurada:

\begin{itemize}
    \item \textbf{Biógrafo incremental} (preparación del corpus): recibe los dominios establecidos y los resúmenes de un periodo, y devuelve la lista acumulativa de dominios y la biografía refinada.
    \item \textbf{Extractor de términos} (consulta, estrategias léxica e híbrida): recibe la idea del estudiante y devuelve dos listas de términos en inglés (generales y específicos).
    \item \textbf{Enriquecedor de consulta} (consulta, estrategias semántica enriquecida e híbrida): recibe la idea y devuelve una reformulación técnica en inglés, con la restricción de no introducir temas nuevos.
\end{itemize}

\pendiente{Volcar aquí el texto literal de los tres prompts (system + plantilla de usuario) y los esquemas Pydantic de sus salidas, una vez congelada la versión final. Incluir un ejemplo de entrada/salida por agente.}

\chapter{Guía de despliegue reproducible}
\label{annex:despliegue}

El sistema completo se despliega con contenedores (Docker Compose) sobre una única máquina anfitriona, coordinados por un proxy inverso que expone todo bajo una única ruta base. La Tabla~\ref{tab:despliegue-servicios} resume los servicios.

\begin{table}[htbp]
    \centering
    \caption{Servicios del despliegue y su cometido.}
    \label{tab:despliegue-servicios}
    \begin{tabularx}{\textwidth}{l X}
        \toprule
        \rowcolor{upmblue!8} \colhead{Servicio} & \colhead{Cometido} \\
        \midrule
        Nginx & Proxy inverso; único punto de entrada externo (enrutado por subruta). \\
        MinIO & Almacén de objetos (capa cruda inmutable y artefactos del pipeline). \\
        Memgraph (+ Lab) & Base de datos de grafo y su interfaz gráfica de exploración. \\
        Qdrant & Índice vectorial (colecciones dispersa BM42 y densa). \\
        Ollama & Servidor de modelos local (LLM y embeddings). \\
        Airflow (+ PostgreSQL) & Orquestación de las fases del pipeline y su base de metadatos. \\
        Aplicación web (Flask) & Interfaz de consulta y recomendación. \\
        \bottomrule
    \end{tabularx}
\end{table}

Los modelos empleados son \texttt{qwen2.5:\allowbreak 7b-instruct} para la generación (biografías y agentes de consulta), \texttt{nomic-\allowbreak embed-\allowbreak text-\allowbreak v2-\allowbreak moe} para los \emph{embeddings} densos y \texttt{Qdrant/\allowbreak bm42-\allowbreak all-\allowbreak minilm-\allowbreak l6-\allowbreak v2-\allowbreak attentions} (BM42) para la representación dispersa.
El desarrollo se llevó a cabo en dos entornos: una estación local (Windows con Docker Desktop) durante las primeras fases y un servidor Linux de la \gls{upm} para el despliegue definitivo (el que sirve la aplicación pública bajo \texttt{wiig.dia.fi.upm.es}), sobre el que opera actualmente todo el sistema.
La Tabla~\ref{tab:hardware} resume las características de ambos.
El contraste es relevante porque demuestra que el diseño es operable tanto en una máquina de desarrollo modesta como en hardware de servidor, y que la holgura de recursos del servidor (en particular, sus dos \glspl{gpu} y su amplia memoria) permite acelerar el enriquecimiento masivo con el \gls{llm} local sin recurrir a servicios externos.

\begin{table}[htbp]
    \centering
    \caption{Características de hardware y software de los dos entornos empleados.}
    \label{tab:hardware}
    \begin{tabularx}{\textwidth}{l X X}
        \toprule
        \rowcolor{upmblue!8} \colhead{Recurso} & \colhead{Desarrollo (local)} & \colhead{Despliegue (servidor)} \\
        \midrule
        Sistema & Windows con Docker Desktop 29.5.3 & Ubuntu 24.04.4 LTS con Docker 29.1.3 \\
        \gls{cpu} & AMD Ryzen 7 7840HS (8 núcleos / 16 hilos) & 2 $\times$ Intel Xeon Gold 5418Y (48 núcleos / 96 hilos) \\
        Memoria & 16\,GB & 503\,GB \\
        \gls{gpu} & NVIDIA GeForce RTX 4060 Laptop (8\,GB) & 2 $\times$ NVIDIA L40 (48\,GB c/u) \\
        Almacenamiento & 1TB & 5,2\,TB \\
        \bottomrule
    \end{tabularx}
\end{table}

\pendiente{Completar al cierre del desarrollo: pasos de arranque desde cero (levantar servicios, ejecutar el pipeline por fases, verificar los índices) y los tiempos de ejecución observados por fase en el servidor.}

El código fuente completo, incluida esta memoria, está disponible en el repositorio público del proyecto: \url{https://github.com/cbadenes/upm-match}.
\pendiente{Si procede, añadir DOI o release etiquetada para la versión entregada.}

\chapter{Resultados extendidos}
\label{annex:resultados-extendidos}

Este anexo desglosa los resultados del banco de pruebas (Sección~\ref{sec:res-benchmark}) por las tres dimensiones que registra el arnés de evaluación: tipo de trabajo, escuela y curso académico.
Para mantener la legibilidad, las tablas reportan la \gls{mrr} de cada estrategia (la métrica que resume la calidad del posicionamiento).
El resto de métricas siguen las mismas tendencias.
En cada fila se resalta en negrita la estrategia con mejor \gls{mrr}.

\begin{table}[htbp]
    \centering
    \caption{\acrshort{mrr} por tipo de trabajo y estrategia.}
    \label{tab:res-ext-tipo}
    \begin{tabularx}{\textwidth}{X c c c c c}
        \toprule
        \rowcolor{upmblue!8} \colhead{Tipo} & \colhead{Léx. dir.} & \colhead{Léx.+LLM} & \colhead{Sem. dir.} & \colhead{Sem. enr.} & \colhead{Híbr.} \\
        \midrule
        \gls{tfg} (227) & 0,049 & 0,090 & 0,145 & \textbf{0,154} & \textbf{0,154} \\
        \gls{tfm} (231) & 0,077 & 0,093 & 0,128 & 0,139 & \textbf{0,153} \\
        \bottomrule
    \end{tabularx}
\end{table}

\begin{table}[htbp]
    \centering
    \caption{\acrshort{mrr} por escuela y estrategia.}
    \label{tab:res-ext-escuela}
    \begin{tabularx}{\textwidth}{X c c c c c}
        \toprule
        \rowcolor{upmblue!8} \colhead{Escuela} & \colhead{Léx. dir.} & \colhead{Léx.+LLM} & \colhead{Sem. dir.} & \colhead{Sem. enr.} & \colhead{Híbr.} \\
        \midrule
        \gls{etsii} (152) & 0,047 & 0,065 & 0,085 & 0,078 & \textbf{0,098} \\
        \gls{etsisi} (97) & 0,046 & 0,075 & 0,147 & 0,153 & \textbf{0,172} \\
        \gls{etsist} (73) & 0,067 & 0,107 & 0,139 & \textbf{0,194} & 0,161 \\
        \gls{etsit} (136) & 0,090 & 0,123 & 0,186 & 0,192 & \textbf{0,198} \\
        \bottomrule
    \end{tabularx}
\end{table}

\begin{table}[htbp]
    \centering
    \caption[\acrshort{mrr} por curso académico y estrategia]{\acrshort{mrr} por curso académico y estrategia. Se omiten los cursos anteriores a 2019 (8 consultas en total, con $n \le 2$ por curso), cuyos valores no son representativos.}
    \label{tab:res-ext-anio}
    \begin{tabularx}{\textwidth}{X c c c c c}
        \toprule
        \rowcolor{upmblue!8} \colhead{Curso} & \colhead{Léx. dir.} & \colhead{Léx.+LLM} & \colhead{Sem. dir.} & \colhead{Sem. enr.} & \colhead{Híbr.} \\
        \midrule
        2019 (14) & 0,005 & 0,149 & \textbf{0,185} & 0,182 & 0,131 \\
        2020 (47) & 0,017 & 0,090 & 0,103 & 0,104 & \textbf{0,154} \\
        2021 (39) & 0,050 & 0,073 & 0,154 & \textbf{0,201} & 0,190 \\
        2022 (66) & 0,066 & 0,097 & 0,118 & 0,128 & \textbf{0,135} \\
        2023 (51) & 0,099 & 0,104 & \textbf{0,174} & 0,167 & 0,159 \\
        2024 (84) & 0,055 & 0,072 & 0,180 & \textbf{0,199} & 0,190 \\
        2025 (122) & 0,092 & 0,105 & 0,108 & 0,111 & \textbf{0,134} \\
        2026 (27) & 0,023 & 0,048 & 0,131 & 0,121 & \textbf{0,142} \\
        \bottomrule
    \end{tabularx}
\end{table}

Tres lecturas se desprenden del desglose.
La ventaja de la señal semántica sobre la léxica es \emph{transversal}: en todos los segmentos la mejor estrategia es siempre una de las tres densas, casi siempre la híbrida o la enriquecida.
Por escuela, la \gls{etsit} es la mejor recuperada y la \gls{etsii} la más difícil, una brecha que probablemente refleja la distinta cobertura de perfiles y la separabilidad temática de cada escuela.
Por tipo, las diferencias entre \gls{tfg} y \gls{tfm} son pequeñas; la híbrida las iguala prácticamente (0,154 frente a 0,153).
El desglose por curso, con tamaños de muestra reducidos en los años antiguos, no muestra una tendencia temporal clara más allá de la ya descrita ventaja de las estrategias densas.

\pendiente{Si se ejecutan varias semillas, añadir desviaciones o intervalos. Considerar el desglose por idioma de la consulta y experimentos secundarios (p.~ej., sensibilidad al límite de recuperación). Añadir los casos ilustrativos que no quepan en el Capítulo~\ref{ch:resultados}.}

El código y esta memoria residen en el repositorio público \url{https://github.com/cbadenes/upm-match}.

La aplicación web está desplegada y es accesible públicamente en \url{https://wiig.dia.fi.upm.es/upm-match/}.